\section{Registration timeline and internal review}
\label{app:timeline}
The title refers to the registered architecture comparisons S1 and S2a.
The project changed direction after earlier results; this sequence was not
one prospectively registered global test. Table~\ref{tab:timeline} separates
prespecified decisions, later authorizations, and supplemental analyses.
Local failures do not falsify every mechanism in their broader family.

\begin{table}[ht]
\centering\small
\begin{tabular}{p{.12\linewidth}p{.24\linewidth}p{.53\linewidth}}
\toprule
Date (2026)&Stage&Disposition and timing qualification\\
\midrule
Oct 21&Finite margin tilts&A2 adds no bound gain over the registered comparator; stopped. Synthetic mechanism test, not a universal negative.\\
Oct 22--25&Multilayer routing&Full request queue completed. Many-history attribution rule added after 189 calls; routing-merger direction stopped under that supplemental rule.\\
Oct 26--27&Gate diagnostics&Partial diagnostic cohort judged gray; PI separately authorized the complete-request cascade test. Gray was not relabeled pass.\\
Oct 27--29&Cascade and local refinements&Main cascade +2 pp [0,4] misses 5 pp. Revised guard test adds one closure, misses five; conditional last-margin test not run.\\
Oct 29--31&S1&Ordinary-training capacity screen fails under both conventions. Secondary incomplete advantage does not rescue it.\\
Oct 31--Nov 1&S2a&Separately authorized direct certification-training test; three primary gates fail. No S2b.\\
Nov 1&Closeout&Fixed two-model reproduction and descriptive analysis authorized after STOP; no new scientific acceptance gate.\\
\bottomrule
\end{tabular}
\caption{Development history. Immutable configurations and decision records
are included in the local artifact. Prior exposure and supplemental timing
prevent treating the entire history as independent confirmation.}
\label{tab:timeline}
\end{table}

The supplied final technical review was written by Claude, an AI also involved
in proposing research directions and gates. It is internal and non-independent.
Codex performed implementation, analysis, and manuscript revisions. Neither
that review nor the environment recreation below is an independent human
review. In revising the manuscript we retained the numerical findings but
removed unsupported priority and causal claims: clean accuracy does not
explain the entire M--D-act deficit, interval width is scale dependent, and
clean top-2 selection does not determine robust-training data exposure.
Submission and human authorship decisions remain with the PI.

\section{Fixed-subset environment reproduction}
\label{app:reproduction}
The PI authorized a final reproduction on D-act seed0 and M seed0, indices
0--49 at $2/255$, with IBP and complete verification once each. It uses fresh,
separately installed environments on the original host, with frozen weights,
materialized input bytes, source identities, and dependencies. Original Python
3.12.12 / torch 2.9.1+cu128 is retained for CPU IBP and Python 3.11.16 /
torch 2.11.0+cpu for CPU verification. Base interpreters and the host are
shared; this is environment reconstruction, not a new-machine or independent
human replication. The numerical implementation and complete-verifier settings
are unchanged. IBP evaluates the requested 50-input subset in one batch per
model rather than the original batch of 256; all resulting numerical
array differences are reported, without changing positive acceptance.

All 200 model--input--mode measurements completed: two 50-input IBP batches
and 100 complete requests. All 100 IBP predictions, labels, clean/certified
flags, and route selections agree with the original records. Primary positives
remain 19/50 for D-act and 12/50 for M. Bounds are not bitwise identical:
maximum absolute lower/upper differences are $9.54\times10^{-6}$ for D-act
and $1.24\times10^{-5}$/$1.43\times10^{-5}$ for M. The changed batch size
is a possible contributor, not an isolated explanation.

Complete terminals agree on 97/100 inputs. D-act agrees on all 50
(21 positive, 22 replayed unsafe, 7 unknown). M retains all 15 positives and
27 replayed unsafe outcomes; its unknown/timeout counts change from 5/3 to
6/2. Specifically, indices 13 and 39 change from TIMEOUT to UNKNOWN, and
index 21 from UNKNOWN to TIMEOUT. These are differences between two unfinished
outcomes, not new certificates or counterexamples. No tolerance was used to
relabel them and no differing input was rerun. The two queue segments took
26.50 minutes of active wall time (85.25 summed request minutes); setup and
saved-record comparison are separate. No execution error or missing request
was recorded. A source-bound record audit and a standalone comparator agree.
These checks do not independently prove the floating-point bounds.

The source snapshot includes nine disclosed textual compatibility edits
across eight files: five original GenBaB edits, a root-input branching-score
edit, and three LiRPA edits. The archive provides their complete diffs and
upstream commits, dependency wheel hashes, weight hashes, fixed input indices,
a ten-call prefix smoke path, and source-based table reconstruction. The
original scientific outcomes are not overwritten by reproduced outcomes.
